% Pursuing Stacks (cotes 134-2 to 134-8): archive cote 134-2, site page 22 (PDF page 23).
% letter to Quillen. Draft compared: ps1.tex:373-433.
% AI-transcribed by Claude Opus 5.5 (claude-opus-5-5, adaptive thinking, effort high), Message Batches API, batch ps1 (msgbatch_01Rv5hCaNTtUt4JHx1yqoHPL), 2026-11-12; not yet checked by eye.
\scanpage{134-2/22}
\mspage{5}
\hand{\unsure{8}}

al structure of a\xout{\unsure{n}} would-be $\infty$- groupoid,  endowed only with its"structure \inserted{of}
order $\leqslant n$\hand{\unsure{"}}. The idea is to give a\xout{\unsure{n}} direct inductive construction of this
sequence,  by describing $C_o$\hand{,} and the passage from $C_n$ to $C_{n+1}$ $(n \geqslant 0)$,  namely
from n-ary to (n+1)-ary structure.  As for the meaning of "universal struc-%
ture", \xout{it depends} once a given structure species is at hand,  it depends
on the type of categories (described by the exactness properties one is assu%
m ing for these) one wants to take as carriers for the considered structure,
and the type of exactness properties one assumes for the functor\hand{s} one allows
between these.  The choice depends partly on the particular species;\xout{xx}
if it is an algebraic structure which can be described say by a handful of
composition laws between a bunch of base sets (or base-\uline{objects}, when loo-%
king at "realizations" of the structure not only in the category (Sets)),
one natural choice is to take categories with finite products, and functors
which commute \deleted{to}\inserted{with} these.  For more sophisticated algebraic structures (inclu-%
ding structure \deleted{of} \hand{"}category\hand{,"} \hand{"}groupoid\hand{"} or the like),  which requires for the
description of data or axioms not only finite products, but also some fiber
products,  one other familiar choice is to take categories with finite in-%
verse limits, and left exact functors.  \over{\illegible{1}}{Still} more sophisticated \deleted{choices}\inserted{structures} ,
when the description of the structure in terms of base objects requires
not only some kind\hand{s} of inverse limits,  but also more or less arbitrary
direct limits (such as the structure of a comodule over an algebra, which
\margin{\unsure{(8) torsors under a group}}
requires consideration of tensor products over a ring object\inserted{(8)}), \deleted{still \xout{\illegible{1}}} \inserted{correspondingly}
more stringent \inserted{exactness} conditions will have to be imposed upon categories and
functors between these,  for the structure to make a sense in \xout{\illegible{1}} these
categories,  and the functors to transform a structure of this type in one
category into one of same type in another.  In most examples I have looked
up,  everything is OK taking categories which are topoi, and functors bet-%
ween these which are inverse image functors for morphisms of topoi, namely
which are left exact and commute with arbitrary direct limits.  There is a
general theorem for the existence of universal structures\inserted{(9)},covering all
these cases  -  for instance there is a "classifying topos" for most \deleted{algebra}\inserted{algebro-}
geometric structures,  whose cohomology say should be viewed as the "classi-%
fying cohomology" of the structure species considered.  In the case we are
interested in here,  it is convenient however to work with the smallest
categories $C_n$ feasable  -  which amounts to being as generous as possible
for the categories \xout{\illegible{2}}  one is allowing  as carriers for the
structure of an $\infty$-groupoid,  and for the functors between these which are
expected to carry an $\infty$-groupoid into an $\infty$-groupoid.  What we will do is
define ultimately a structure of an $\infty$-groupoid in a category C,  as a
sequence of objects $F_i$ $(i\hand{\in} \underline{N})$ \hand{\unsure{,}} endowed with some structure to be defined,

\margin{22}
